% Section 6, part 9.
The base case $h=1$ of Theorem~V rests on a local formula for $\omega_{1,1}$. Let $c$ be a branch point, $t:=p-c$, and $\chi$ (a local letter) a local coordinate at $c$ with $x=x(c)+\lambda\chi^2$, $\lambda\neq 0$ a constant, so that $\iota_c\colon\chi\mapsto-\chi$. Write $y=\sum_ky_k\chi^k$ and, near $(c,c)$, $\omega_{0,2}=\big[(\chi_1-\chi_2)^{-2}+\mathrm{reg}(\chi_1,\chi_2)\big]\dd\chi_1\dd\chi_2$ with $\mathrm{reg}$ holomorphic and symmetric; then $\mathrm{reg}_{00}:=\mathrm{reg}(0,0)=\frac16\{t;\chi\}(0)$, where $\{t;\chi\}=t'''/t'-\frac32(t''/t')^2$ is the Schwarzian derivative of $t$ as a function of $\chi$ (the regular part on the diagonal of $t'(\chi_1)t'(\chi_2)/(t(\chi_1)-t(\chi_2))^2$). For a $1$-form $\eta=\sum_m\eta_m\chi^m\dd\chi$ holomorphic at $c$ put
\begin{equation}
\label{eq:mod-loc}
\mathrm{Loc}_c[\eta]:=-\frac{1}{4\lambda y_1}\Big[\mathrm{reg}_{00}\,\eta_0+\frac{\eta_2}{12}-\frac{y_3}{4y_1}\,\eta_0\Big].
\end{equation}

\begin{lem}[A local formula for $\omega_{1,1}$]
\label{lem:mod-w11}
\begin{enumerate}
\item $\omega_{1,1}(p_0)=\sum_c\mathrm{Loc}_c[\omega_{0,2}(p_0,\cdot)]$ for $p_0$ off the branch points, $\omega_{0,2}(p_0,\cdot)$ being a form in the second slot.
\item $\omega_{1,1}=\dd G$ with $G$ rational, with poles only at the branch points and $G(\infty)=0$; so $\omega_{1,1}$ has zero residues, is holomorphic at $p=0$ and is $O(\dd p/p^2)$ at $\infty$.
\item $\langle-\log p\rangle^{\mathrm{TR}}_1=\int_{\mathfrak r}\omega_{1,1}=\sum_c\mathrm{Loc}_c[-\dd q/q]$.
\end{enumerate}
\end{lem}

\begin{proof}
(i) By \eqref{eq:mod-kernel},
\[
\omega_{1,1}(p_0)=\sum_c\Res_{q\to c}K_c(p_0,q)\,\omega_{0,2}(q,\iota_cq).
\]
Expand $\omega_{0,2}(p_0,q')=\dd p_0\sum_mB_m(p_0)\chi^m\dd\chi$ for $q'$ near $c$ with coordinate $\chi$. Then $\int_{\iota_cq}^q\omega_{0,2}(p_0,\cdot)=2\,\dd p_0\sum_{m\text{ even}}B_m\chi^{m+1}/(m+1)$, $y(q)-y(\iota_cq)=2\sum_{k\text{ odd}}y_k\chi^k$ and $\dd x=2\lambda\chi\,\dd\chi$, so
\[
K_c(p_0,q)=\frac{\dd p_0}{\dd\chi}\,\frac{1}{4\lambda y_1}\Big[B_0\chi^{-1}+\Big(\frac{B_2}{3}-B_0\frac{y_3}{y_1}\Big)\chi+O(\chi^3)\Big],
\]
and $\omega_{0,2}(q,\iota_cq)=-\big[(4\chi^2)^{-1}+\mathrm{reg}_{00}+O(\chi^2)\big]\dd\chi^2$, since $\dd(-\chi)=-\dd\chi$ and $\mathrm{reg}(\chi,-\chi)=\mathrm{reg}_{00}+O(\chi^2)$ by symmetry. The coefficient of $\chi^{-1}\dd\chi$ in the product is
\[
-\frac{1}{4\lambda y_1}\Big[B_0\,\mathrm{reg}_{00}+\frac14\Big(\frac{B_2}{3}-B_0\frac{y_3}{y_1}\Big)\Big]\dd p_0=\mathrm{Loc}_c[\omega_{0,2}(p_0,\cdot)].
\]
(ii) $\omega_{0,2}(p_0,q)=\dd_{p_0}\big[\dd q/(q-p_0)\big]$, and near $c$, $\dd q/(q-p_0)=\sum_mg_m(p_0)\chi^m\dd\chi$ with $g_m$ rational in $p_0$, with poles only at $c$ and $g_m(\infty)=0$ (expand $1/(c+t(\chi)-p_0)$ in powers of $t$). So $B_m(p_0)\dd p_0=\dd g_m(p_0)$, and by (i) $\omega_{1,1}$ is the differential of a linear combination $G$ of the $g_0,g_2$ at the five branch points. (iii) Each residue in (i) is an integral over a small circle around $c$, which $\mathfrak r$ avoids; by (ii) the ray integral converges, and since $\omega_{0,2}(p_0,q)=O(p_0^{-2})\dd p_0$ uniformly for $q$ on the circles, the integral over $p_0\in\mathfrak r$ may be taken inside: $\int_{p_0\in\mathfrak r}\omega_{0,2}(p_0,q)=\dd q\big[-1/(p_0-q)\big]_{p_0=0}^{p_0=\infty}=-\dd q/q$ for $q$ off the ray, and $\mathrm{Loc}_c$ is linear in finitely many Taylor coefficients.
\end{proof}

By (i) and \eqref{eq:mod-kernel}, $\mathrm{Loc}_c[\eta]=\Res_{q\to c}\frac12\frac{\int_{\iota_cq}^q\eta}{(y(q)-y(\iota_cq))\dd x(q)}\,\omega_{0,2}(q,\iota_cq)$ for $\eta$ holomorphic at $c$; this form does not depend on the coordinate $\chi$.

\begin{thm}[Theorem V]
\label{thm:mod-V}
Let $h\geq 0$, $a,k\geq 0$, $n=a+k\geq 1$, $2h-2+n>0$, and $\phi_1,\dots,\phi_a\in\bbQ[v][p]$.
\begin{enumerate}
\item If $a\geq 1$, then $\langle\phi_1\cdots\phi_a,(-\log p)^k\rangle^{\mathrm{TR}}_h\in\bbQ[v]$.
\item If $a=0$, then $\langle(-\log p)^k\rangle^{\mathrm{TR}}_h=(\FTR_h)^{(k)}(v)$, with $\FTR_h$ of \cref{thm:mod-gap} for $h\geq 2$, $\FTR_1:=-\frac1{12}\log v$ (here $k\geq 1$) and $\FTR_0:=\frac12v^2\log v$ (here $k\geq 3$).
\item The recursion \eqref{eq:mod-V1} holds.
\end{enumerate}
\end{thm}

The quoted inputs of the proof are \cite[Thm.~5.1, (5-14)]{Eynard:2007kz} (through \cref{thm:mod-var}~(i) and \cref{prop:mod-V1}) and \cite[Thm.~7.1]{ABDKS25} (through \cref{thm:mod-Lreg} and \cref{thm:mod-gap}), besides \cite[Thm.~4.2, Thm.~4.6]{Eynard:2007kz} (through the paragraph after \cref{conv:mod-kernel}, \cref{def:mod-pair} and \cref{lem:mod-xfun}).

\begin{proof}
(iii) is \cref{prop:mod-V1}. (i) Fix $h$ and $a\geq 1$; we use induction on $k$. If $2h-2+a>0$, the case $k=0$ is \cref{thm:mod-Lreg}. If $2h-2+a\leq 0$, then $(h,a)$ is $(0,1)$ or $(0,2)$, the first admissible pairing is $\langle\phi_1,(-\log p)^2\rangle^{\mathrm{TR}}_0$, respectively $\langle\phi_1\phi_2,-\log p\rangle^{\mathrm{TR}}_0$, and it lies in $\bbQ[v]$ by \cref{lem:mod-genus0}~(iv). For the step from $k$ to $k+1$ (with $2h-2+a+k>0$), \eqref{eq:mod-V1} writes the pairing with $k+1$ log legs through the $v$-derivative of a pairing with $k$ log legs and pairings with the same number $a$ of L-legs and $k$ log legs, the legs $\partial_v\phi_i$ and $(\cX\dd\phi_i)_+$ lying in $\bbQ[v][p]$ because $p^5-v$ is monic; all of these lie in $\bbQ[v]$ by the induction hypothesis. The L-legs are never consumed, so the induction stays in the class $a\geq 1$.

(ii) For $a=0$, \eqref{eq:mod-V1} reads $\langle(-\log p)^{k+1}\rangle^{\mathrm{TR}}_h=\partial_v\langle(-\log p)^k\rangle^{\mathrm{TR}}_h$, so it suffices to treat the first admissible $k$. \emph{$h\geq 2$, $k=1$.} As in the proof of \cref{prop:mod-period}~(iii), $(2-2h)\FTR_h=\frac15\langle p^5\rangle^{\mathrm{TR}}_h+v\langle-\log p\rangle^{\mathrm{TR}}_h$ with $\langle p^5\rangle^{\mathrm{TR}}_h=0$; since $\FTR_h$ is a multiple of $v^{2-2h}$ (\cref{thm:mod-gap}), $\langle-\log p\rangle^{\mathrm{TR}}_h=(2-2h)\FTR_h/v=(\FTR_h)'(v)$. This case uses neither \cref{thm:mod-var} nor anything beyond \cref{thm:mod-Lreg} and \cref{thm:mod-gap}. \emph{$h=0$, $k=3$.} $\langle(-\log p)^3\rangle^{\mathrm{TR}}_0=1/v=(\FTR_0)'''$ by \cref{lem:mod-genus0}~(iv). \emph{$h=1$, $k=1$.} We show $\langle-\log p\rangle^{\mathrm{TR}}_1=-1/(12v)$.

\emph{Reduction to $v=1$.} By \cref{lem:mod-scale}~(ii) with $k=\varsigma$, $m_\varsigma^*\omega^{(v)}_{1,1}=\varsigma^{5(2-2-1)}\omega^{(1)}_{1,1}=v^{-1}\omega^{(1)}_{1,1}$; $m_\varsigma$ maps rays to rays, and the pairing does not depend on the ray. So $\langle-\log p\rangle^{\mathrm{TR}}_1(v)=v^{-1}\langle-\log p\rangle^{\mathrm{TR}}_1(1)$. Let $v=1$, $\varsigma=1$, so $\varsigma_b=\zeta^{-b}$.

\emph{Five equal contributions.} The map $\rho_b(p):=\zeta^{-b}p$ sends $1$ to $\varsigma_b$ and satisfies $x\circ\rho_b=\zeta^bx$, $y\circ\rho_b=\zeta^{-b}y$, $\rho_b^*\omega_{0,2}=\omega_{0,2}$ and $\rho_b^*(\dd q/q)=\dd q/q$; hence $\iota_b\circ\rho_b=\rho_b\circ\iota_0$, and the factor $(y(q)-y(\iota q))\,\dd x(q)$ is unchanged ($\zeta^{-b}\zeta^b=1$). By the coordinate-free form of $\mathrm{Loc}$ after \cref{lem:mod-w11}, $\mathrm{Loc}_{\varsigma_b}[-\dd q/q]=\mathrm{Loc}_1[-\dd q/q]$ for every $b$.

\emph{The contribution at $c=1$.} With $t=p-1$, $x(1+t)=\frac54+\frac52t^2+0\cdot t^3+\frac54t^4-t^5+O(t^6)$, so $\lambda=5/2$ and $\chi=t\big(1+\frac{t^2}2-\frac{2t^3}5+O(t^4)\big)^{1/2}=t+\frac{t^3}4-\frac{t^4}5+O(t^5)$, that is $t=\chi-\frac{\chi^3}4+\frac{\chi^4}5+O(\chi^5)$. Hence $y=1+t$ has $y_1=1$, $y_2=0$, $y_3=-\frac14$; $t'(0)=1$, $t''(0)=0$, $t'''(0)=-\frac32$, so $\{t;\chi\}(0)=-\frac32$ and $\mathrm{reg}_{00}=-\frac14$; and $-\dd q/q=-t'(\chi)\,\dd\chi/(1+t(\chi))=-\big(1-\chi+\frac14\chi^2+O(\chi^3)\big)\dd\chi$, so $\eta_0=-1$, $\eta_2=-\frac14$. By \eqref{eq:mod-loc},
\[
\mathrm{Loc}_1[-\dd q/q]=-\frac1{10}\Big[\Big(-\frac14\Big)(-1)+\frac{-1/4}{12}-\frac{-1/4}{4}\,(-1)\Big]=-\frac1{10}\cdot\frac{12-1-3}{48}=-\frac1{60}.
\]
By \cref{lem:mod-w11}~(iii) the total at $v=1$ is $5\cdot(-\frac1{60})=-\frac1{12}$. So $\langle-\log p\rangle^{\mathrm{TR}}_1=-1/(12v)=(\FTR_1)'(v)$, which also fixes $\FTR_1$ up to an additive constant.
\end{proof}

% Section 6, part 10.
\subsection{The Model Lemma on \texorpdfstring{$\sK$}{K}}
\label{sec:mod-MLK}

Recall \cref{def:red-ML}: $\ML{h}$ asserts (i) that for $n\geq 1$, $2h-2+n>0$ and $X_1,\dots,X_n\in\sK$ homogeneous over $\bbQ[v]$ (that is $X_i\in\sK\cap V_{w_i}$), $\langle X_1\cdots X_n\rangle_h(v)-(\Fm_h)^{(n)}(v)\prod_i\Pi_0(X_i)$ lies in $\bbQ[v]$, and (ii) that $\Fm_h=c_hv^{2-2h}$ for $h\geq 2$; here $\Fm_h$ is the free energy of \cref{cor:mod-trF} for $h\geq 2$, and $\Fm_1=-\frac1{12}\log v$, $\Fm_0=-\frac52v^2\log v$, of which only the derivatives of order $\geq 1$, respectively $\geq 3$, enter. Note that $\Fm_h=(-5)^{1-h}\FTR_h$ for every $h\geq 0$, with $\FTR_0,\FTR_1$ as in \cref{thm:mod-V}~(ii).

\begin{thm}[The Model Lemma on $\sK$]
\label{thm:mod-MLK}
$\ML{h}$ holds for every $h\geq 0$.
\end{thm}

\begin{proof}
Part (ii) of $\ML{h}$ is \cref{thm:mod-gap}. For part (i), let $X_i\in\sK\cap V_{w_i}$, $i=1,\dots,n$.

\emph{Step 1 (topological recursion).} By \cref{thm:mod-tr}, $\langle q_{\phi_1}\cdots q_{\phi_n}\rangle_h=(-5)^{1-h}\langle\phi_1\cdots\phi_n\rangle^{\mathrm{TR}}_h$ for legs as in \cref{def:mod-ancvec}.

\emph{Step 2 (the leg dictionary).} By \cref{prop:mod-LD}, the TR leg of $Y\in V_w$ is $-\Lambda(Y)$, and $q_{-\log p}=H$.

\emph{Step 3 (the legs of $\sK$).} By \cref{cor:mod-Ksplit}, $X_i=Y_i+\Pi_0(X_i)H$ with $Y_i\in\sKo\cap V_{w_i}$ and $\Pi_0(X_i)\in\bbQ v^{w_i/5}$, a monomial with a non-negative exponent. Hence $X_i=q_{\ell_i}$ for the leg $\ell_i:=-\Lambda(Y_i)+\Pi_0(X_i)\,(-\log p)$ ($q$ is $\bbQ[v]$-linear in the leg). By \cref{thm:mod-lattice}~(b) and \cref{lem:mod-classL}~(i), $-\Lambda(Y_i)=\phi_i+f_i(x)$ with $\phi_i\in\bbQ[v][p]$ and $f_i\in\bbQ[v][x]$. (It is the lattice theorem, not homogeneity, that excludes terms $v^{-j}p^{w_i+5j}$ with $j>0$, which homogeneity would allow.)

\emph{Step 4 (expansion).} By multilinearity and \cref{lem:mod-xfun}, which removes the components $f_i(x)$ (in every slot, whatever the other slots carry),
\[
\langle X_1\cdots X_n\rangle_h=(-5)^{1-h}\sum_{S\subseteq\{1,\dots,n\}}\prod_{i\notin S}\Pi_0(X_i)\,\big\langle(\phi_i)_{i\in S},(-\log p)^{n-|S|}\big\rangle^{\mathrm{TR}}_h .
\]
Every term with $S\neq\emptyset$ lies in $\bbQ[v]$: by \cref{thm:mod-Lreg} if $S=\{1,\dots,n\}$, and by \cref{thm:mod-V}~(i) otherwise; the factors $\Pi_0(X_i)$ are monomials with non-negative exponents.

\emph{Step 5 (the pure logarithmic term).} By \cref{thm:mod-V}~(ii) and \cref{cor:mod-trF} (for $h=0,1$ by the normalisations recalled above), the term $S=\emptyset$ is
\[
(-5)^{1-h}\langle(-\log p)^n\rangle^{\mathrm{TR}}_h\prod_i\Pi_0(X_i)=(-5)^{1-h}(\FTR_h)^{(n)}\prod_i\Pi_0(X_i)=(\Fm_h)^{(n)}\prod_i\Pi_0(X_i).
\]
This proves (i).
\end{proof}

\Cref{thm:mod-MLK} holds modulo the quoted inputs named at the beginning of this section, and nothing about the quintic enters. The insertions met in \cref{thm:red-main} are $\bbQ\fps{u}$-combinations of homogeneous elements of $\sK$ (\cref{lem:tra-coeff}); there the statement is applied termwise (\cref{lem:tra-conv}~(b)).

\begin{cor}[The one-point law on $\sK$]
\label{cor:mod-onepoint}
Let $h\geq 1$ and $X\in\sK\cap V_w$. If $\langle X\rangle_h$ is singular, that is, if its $v$-exponent is negative, then $\langle X\rangle_h=\Pi_0(X)\,(\Fm_h)'(v)$.
\end{cor}

\begin{proof}
By \cref{thm:mod-MLK} with $n=1$, $\langle X\rangle_h-(\Fm_h)'\Pi_0(X)\in\bbQ[v]$. By \cref{lem:set-modhom}, $\langle X\rangle_h=Cv^e$ with $e=1-2h+w/5$, and $(\Fm_h)'\Pi_0(X)$ is a multiple of the same power $v^{1-2h}\cdot v^{w/5}$, since $\Pi_0(X)\in\bbQ v^{w/5}$; as $e<0$ by assumption, the difference, an element of $\bbQ[v]$ of this form, is $0$.
\end{proof}

\begin{cor}[The primary sector]
\label{cor:mod-primary}
Let $h\geq 1$ and $n\geq 1$. A primary correlator $\langle H^{a_1}\cdots H^{a_n}\rangle_h$ vanishes unless $a_1=\dots=a_n=1$, and $\langle H^n\rangle_h=(\Fm_h)^{(n)}(v)$.
\end{cor}

\begin{proof}
By \cref{def:tra-tower} and the shift rule, $\Pi_i(H^a)=0$ for $i\geq 1$, and $\Pi_0(H^a)=0$ for $a\neq 1$, $\Pi_0(H)=1$; so $H^a\in\sKo$ for $a\neq 1$ and $H\in\sK$. By \cref{lem:set-modhom} the $v$-exponent of a primary correlator is $2-2h+\sum_i\wt(0,a_i)$ with $\wt(0,a)=(\nu_a-4)/5<0$, negative for $h\geq 1$. By \cref{thm:mod-MLK} the difference $\langle\prod_iH^{a_i}\rangle_h-(\Fm_h)^{(n)}\prod_i\Pi_0(H^{a_i})$ lies in $\bbQ[v]$ and is homogeneous of negative exponent, hence $0$; and $\prod_i\Pi_0(H^{a_i})$ is $1$ if all $a_i=1$ and $0$ otherwise.
\end{proof}

\begin{rmk}[Off $\sK$]
\label{rmk:mod-offK}
The restriction to $\sK$ cannot be dropped. Off $\sK$, that is, for insertions with a component along the directions $\psi^kH$ with $k\geq 1$, the singular part of a one-point correlator is in general not $\Pi_0(X)(\Fm_h)'$; and the singular parts of two-point correlators off $\sK$ are not of the form given by $\Fm_h(V(q))$ for any function $V$ of the insertion independent of the genus, linear or not. Both facts are known only by computation (genus $3$ and $4$ examples; \cref{sec:evi-model}); they are not used.
\end{rmk}

